
The gradient of $f$ is
\[
\nabla f(x) = \cvect{\frac{\partial f}{\partial x_1} \\ \vdots
  \\ \frac{\partial f}{\partial x_n}} = \cvect{c_1 \\ \vdots \\ c_n} = c,
\]
for any $x \in \R ^n$.
As it is a constant vector, the second derivative matrix is zero, for
any $x \in \R^n$:
\[
\nabla^2 f(x) = 0.
\]
We first note that the value of $b$ is irrelevant to determine the
necessary optimality conditions. Adding a constant to an objective
function does not change the optima. 
If $c=0$, that is, if
\[
c_1 = \ldots = c_n = 0,
\]
then the necessary optimality conditions are verified for all
$x\in\R^n$.
If $c \neq 0$, that is, if there is at least one $i$ such that $c_i
\neq 0$, then the first order necessary optimality conditions are not
verified . The second order necessary conditions are always verified,
as the null matrix is positive semidefinite.

Note that, if $c=0$, the linear function is constant, with value
$b$. And any $x\in \R^n$ is a (global) optimum. If not, the function
is not constant. As there is no constraint, it is not bounded and
there is no optimum.  

